\section{Especificação da Linguagem}
\label{sec:especificacao}

\subsection{Padrões oficiais}

\textbf{Padrão Mercosul.} Adotado em novos emplacamentos e em
mudanças de município ou estado \cite{contran780}. Sua sequência é
composta por três letras, um número, uma letra e dois números
(LLLNLNN), em que L é letra e N é número \cite{contran741}.

\textbf{Padrão anterior ao Mercosul.} Segue a sequência de
três letras seguidas de quatro números. O hífen presente na
chapa física (por exemplo, \texttt{ABC-1234}) é apenas visual e não
faz parte da sequência de caracteres.

\subsection{Alfabeto}

Sejam $\Sigma_L = \{A, B, \dots, Z\}$ o conjunto das letras
maiúsculas e $\Sigma_N = \{0, 1, \dots, 9\}$ o conjunto dos
algarismos. O alfabeto do autômato é
\[
  \Sigma = \Sigma_L \cup \Sigma_N \cup \{-\}.
\]

\subsection{Decisões de projeto}

\begin{enumerate}
  \item São aceitas apenas as 26 letras do alfabeto latino em caixa
        alta, sem acentos, cedilha ou caracteres especiais.
  \item São aceitos os 10 algarismos arábicos, de 0 a 9.
  \item O hífen é aceito de forma opcional, apenas no padrão anterior
        ao Mercosul, entre a terceira letra e o primeiro número.
  \item Em placas convertidas do sistema antigo, a quinta posição
        assume apenas letras de A a J. Enquanto o autômato Mercosul
        geral (AFD 02) aceita qualquer letra de A a Z nessa posição,
        foi implementado um autômato específico (AFD 03) para validar
        essa restrição de A a J.
  \item O autômato valida o formato da placa, e não a sua existência.
\end{enumerate}

\subsection{Expressão regular}

\[
  \Lt\,\Lt\,\Lt\,\big(\,\Hf?\,\Nm\,\Nm\,\Nm\,\Nm \;\mid\; \Nm\,\Lt\,\Nm\,\Nm\,\big)
\]

em que $\Lt$, $\Nm$ e $\Hf$ representam, respectivamente, uma letra
maiúscula, um algarismo e o hífen.

\subsection{Exemplos iniciais de validação}

\begin{table}[h]
\centering
\begin{tabular}{lllp{6cm}}
\toprule
Entrada & Padrão & Resultado & Motivo \\
\midrule
\texttt{BRA2E19}  & Mercosul & Válida   & Segue LLLNLNN \\
\texttt{QTP9Z01}  & Mercosul & Válida   & Letra de A a Z na 5ª posição \\
\texttt{ABC1234}  & Anterior & Válida   & Segue LLLNNNN \\
\texttt{ABC-1234} & Anterior & Válida   & Hífen opcional (decisão de projeto) \\
\texttt{ABC-1D23} & ---      & Inválida & Hífen no padrão Mercosul \\
\texttt{AB12C34}  & ---      & Inválida & Apenas 2 letras iniciais \\
\texttt{ABC12D3}  & ---      & Inválida & Letra na 6ª posição \\
\texttt{BRA2E1}   & ---      & Inválida & Apenas 6 caracteres \\
\bottomrule
\end{tabular}
\caption{Exemplos iniciais de entradas válidas e inválidas do validador.}
\label{tab:exemplos-especificacao}
\end{table}